\section{Relative quotient action}

Предыдущий gate вывел каноническую derivation
\(\delta_h(F)=\frac12[h,F]\), но оставил выбор её нормы внешним. Этот
выбор можно устранить, если curvature рассматривается как относительный
класс по отношению к fixed-point algebra трёхузловой цепи. Получающийся
функционал является quotient norm, а не вручную выбранным matrix projector.

\subsection{Fixed-point algebra и conditional expectation}

Пусть
\begin{equation}
 h=\operatorname{diag}(-I_4,0_2,I_4)
\end{equation}
и
\begin{equation}
 \alpha_t(X)=e^{ith}Xe^{-ith}.
\end{equation}
Поскольку spectrum \(h\) целочислен, \(\alpha_t\) задаёт каноническое
действие окружности. Его fixed-point algebra равна
\begin{equation}
 \mathcal B_h=\operatorname{Fix}(\alpha)
 =\operatorname{End}(\bar4)\oplus
  \operatorname{End}(2_R)\oplus\operatorname{End}(4).
 \label{eq:ps-relative-fixed-algebra}
\end{equation}
Единственное сохраняющее обычный trace conditional expectation на
\(\mathcal B_h\) имеет вид
\begin{equation}
 E_h(X)=\frac1{2\pi}\int_0^{2\pi}\alpha_t(X)\,dt.
 \label{eq:ps-relative-expectation}
\end{equation}
В конечномерной Hilbert--Schmidt metric это ортогональная projection на
block-diagonal fixed-point algebra.

\subsection{Исключение общей auxiliary curvature}

Определим относительное действие не на представителе \(F\), а на его классе
в Hilbert quotient по \(\mathcal B_h\):
\begin{equation}
 S_{\rm quot}[F]
 =\inf_{C\in\mathcal B_h}\|F-C\|_F^2.
 \label{eq:ps-relative-quotient-action}
\end{equation}
Ортогональность даёт точное разложение
\begin{equation}
 \|F-C\|_F^2
 =\|F-E_h(F)\|_F^2+\|E_h(F)-C\|_F^2.
 \label{eq:ps-relative-pythagoras}
\end{equation}
Следовательно, минимум единственен:
\begin{equation}
 C_\star=E_h(F),
 \qquad
 \boxed{S_{\rm quot}[F]=\|F-E_h(F)\|_F^2.}
 \label{eq:ps-relative-auxiliary-elimination}
\end{equation}
Тем самым relative norm возникает после algebraic elimination общей
block-diagonal curvature. Непрерывный coefficient или дополнительный
projector внутри quotient construction не вводится. Общий вес этого сектора
относительно исходного Pati--Salam trace рассматривается отдельно ниже.

\subsection{Совпадение с mapping-cone Dirichlet form}

Для \(F=\mathcal D_\Delta^2\) parity запрещает nearest-neighbour blocks:
остаются только diagonal backtracking blocks и endpoint blocks
\(0\leftrightarrow2\). На этом чётном curvature space оператор
\begin{equation}
 Q_h=\frac14\operatorname{ad}_h^2
\end{equation}
имеет eigenvalue ноль на \(\mathcal B_h\) и eigenvalue единицу на endpoint
blocks. Поэтому
\begin{align}
 F-E_h(F)&=Q_hF,\\
 S_{\rm quot}[F]
 &=\langle F,Q_hF\rangle_F
 =\left\|\frac12[h,F]\right\|_F^2.
 \label{eq:ps-relative-quotient-dirichlet}
\end{align}
Endpoint gap координаты \((-1,0,1)\) равен двум, поэтому factor \(1/4\)
фиксирован нормировкой graph distance.

Подстановка admissible Pati--Salam curvature даёт
\begin{equation}
 \boxed{
 S_{\rm quot}[\mathcal D_\Delta^2]
 =4\det(\Delta\Delta^\dagger).}
 \label{eq:ps-relative-quotient-selector}
\end{equation}

\begin{theorem}[Canonical relative quotient action]
На чётном curvature space цепи \(\bar4\to2_R\to4\) trace-preserving
conditional expectation \(E_h\) однозначно задаёт quotient norm по
fixed-point algebra. После нормировки endpoint graph distance к двум этот
quotient norm совпадает с \(\|\frac12[h,F]\|_F^2\) и на
\(F=\mathcal D_\Delta^2\) равен
\(4\det(\Delta\Delta^\dagger)\).
\end{theorem}

\subsection{Gauge covariance и граница применимости}

Block-diagonal Pati--Salam gauge transformations коммутируют с \(h\), поэтому
\begin{equation}
 E_h(UFU^\dagger)=UE_h(F)U^\dagger
\end{equation}
и \(S_{\rm quot}\) gauge invariant. Ограничение на even curvature
существенно. На полном matrix space odd nearest-neighbour blocks имеют
height gap единицу; там \(\frac14\operatorname{ad}_h^2\) имеет eigenvalue
\(1/4\), тогда как \(I-E_h\) имеет eigenvalue единицу. Поэтому найденное
тождество нельзя переносить на произвольный operator.

\subsection{Vacuum consequence и общий вес}

Пусть \(\lambda_{\rm rel}\) обозначает общий вес relative carrier в полном
bosonic trace. Тогда composite potential принимает вид
\begin{equation}
 V_\Delta=-\rho^2+\tau^2
 +4\lambda_{\rm rel}\det(\Delta\Delta^\dagger).
\end{equation}
На rank-one Standard-Model orbit exact Hessian spectrum равен
\begin{equation}
 8\sqrt2\ (1),
 \qquad
 0\ (9),
 \qquad
 \sqrt2(4\lambda_{\rm rel}-2)\ (6).
 \label{eq:ps-relative-parent-hessian}
\end{equation}
Следовательно, устойчивость требует
\begin{equation}
 \lambda_{\rm rel}>\frac12.
\end{equation}
Для one-copy embedding в ту же normalized trace--Hodge metric
\(\lambda_{\rm rel}=1\), и последние шесть eigenvalues равны
\(2\sqrt2\). Quotient construction фиксирует внутренний coefficient четыре,
но не доказывает сама по себе multiplicity и общий trace weight carrier.

\begin{theorem}[Условное закрытие parent-norm gate]
Свобода выбора projector и его внутренней normalization закрыта. Полное
физическое замыкание требует вывести auxiliary/common status, multiplicity
и общий trace weight \(\lambda_{\rm rel}\) из полного Pati--Salam relative
spectral triple или BV complex, а затем пересчитать endpoint-state, mixed
Hessian и gauge-running ledger.
\end{theorem}

\subsection{Литература и воспроизводимость}

Relative spectral triples и boundary quotient data сверены с I.~Forsyth,
M.~Goffeng, B.~Mesland and A.~Rennie, arXiv:1607.07143. Tracial conditional
expectation сопоставлена с E.~A.~Carlen and A.~Vershynina,
arXiv:1710.02409. Higher-derivative Yang--Mills terms spectral action
сверены с W.~D.~van Suijlekom, arXiv:1104.5199. Конечномерная identity
\eqref{eq:ps-relative-quotient-selector} является специализацией проекта.

Machine audit сохранён в
\path{s2t_v4_pati_salam_relative_parent_action_gate.py} и
\path{s2t_v4_pati_salam_relative_parent_action_gate_results.json}.